\section*{अध्याय \ref{ch:b3:palaeomagnetism} --- भू-चुंबकत्व और पुराचुंबकत्व}

\begin{solution}{exo:b3:palaeomagnetism:1}
$\tan I = 2\tan\ang{30} = 1.1547$, अतः $I = \ang{49.1}$। इसके विपरीत, $\tan\lambda
= \tfrac12\tan(\ang{-20}) = -0.1820$, अतः $\lambda = \ang{-10.3}$, अर्थात्
\ang{10.3}\,S।
\end{solution}

\begin{solution}{exo:b3:palaeomagnetism:2}
$H = \sqrt{20000^2 + 2000^2} = \qty{20100}{nT}$;
$F = \sqrt{20000^2 + 2000^2 + 40000^2} = \qty{44770}{nT}$। चुंबकीय
क्रांति $\arctan(-2000/20000) = \ang{-5.7}$ है, उत्तर से \ang{5.7} पश्चिम की ओर;
नति $\arctan(40000/20100) = \ang{63.3}$, नीचे की ओर।
\end{solution}

\begin{solution}{exo:b3:palaeomagnetism:3}
$B_0 = -g_1^0 = \qty{29350}{nT}$ भूमध्य रेखा पर, और $2B_0 = \qty{58700}{nT}$
ध्रुवों पर। $m = 4\pi a^3B_0/\mu_0 = a^3B_0/\num{e-7} = \num{2.5862e20}\times
\num{2.935e-5}/\num{e-7} = \qty{7.59e22}{A.m^2}$, जो
$g_1^1$ और $h_1^1$ जोड़ने वाले पूर्ण द्विध्रुव के \qty{7.69e22}{A.m^2} से थोड़ा कम है।
\end{solution}

\begin{solution}{exo:b3:palaeomagnetism:4}
(क) \omterm{def:b3:palaeomagnetism:nrm}{ऊष्मा-अवशिष्ट चुंबकन}, जो लावा के ठंडा होते समय
अपने कणों के अवरुद्धन तापमानों से गुज़रते हुए क्षेत्र अभिलेखित करता है। (ख) \omterm{def:b3:palaeomagnetism:nrm}{अपरद-अवशिष्ट चुंबकन}, जो
गाद को बैठने और
संहनन द्वारा जकड़े जाने में लगे समय के भीतर क्षेत्र अभिलेखित करता है। (ग) \omterm{def:b3:palaeomagnetism:nrm}{रासायनिक अवशिष्ट चुंबकन}, जो
हेमेटाइट की वृद्धि के दौरान का क्षेत्र, प्रारंभिक दफ़न के समय और इसलिए
निक्षेपण के कुछ बाद अभिलेखित करता है। (घ) वर्तमान क्षेत्र के अनुदिश श्यान अवशिष्ट चुंबकन, जो
मूल \omterm{def:b3:palaeomagnetism:nrm}{ऊष्मा-अवशिष्ट चुंबकन} में सबसे कम शिथिलन काल वाले कणों द्वारा जुड़ता है; इसे
विचुंबकन के पहले चरणों में हटा दिया जाता है।
\end{solution}

\begin{solution}{exo:b3:palaeomagnetism:5}
इकाई सदिश $(\cos I\cos D,\ \cos I\sin D,\ \sin I)$ हैं
\begin{gather*}
  (0.6545, 0.1391, 0.7431),\quad (0.6670, 0.2695, 0.6947),\\
  (0.5235, 0.1501, 0.8387),\quad (0.5985, 0.0629, 0.7986).
\end{gather*}
योग $(2.4435, 0.6216, 3.0751)$ है, जिसकी लंबाई $R = 3.9766$ है। माध्य:
$D = \arctan(0.6216/2.4435) = \ang{14.3}$, $I = \arcsin(3.0751/3.9766) =
\ang{50.7}$। $k = 3/(4 - 3.9766) = 128$। $20^{1/3} - 1 = 1.7144$ लेने पर,
$\cos\alpha_{95} = 1 - (0.0234/3.9766)\times1.7144 = 0.98991$, अतः
$\alpha_{95} = \ang{8.1}$ है: चार नमूनों के साथ शंकु चौड़ा ही रहता है, चाहे
समूह कितना भी कसा हो।
\end{solution}

\begin{solution}{exo:b3:palaeomagnetism:6}
$\lambda = \arctan(\tfrac12\tan\ang{50}) = \ang{30.79}$, अतः $p = \ang{59.21}$।
\newline $\sin\lambda_p = \sin\ang{40}\cos\ang{59.21} + \cos\ang{40}\sin\ang{59.21}
\cos\ang{20} = 0.9474$, $\lambda_p = \ang{71.34}$\,N। $\sin\beta =
\sin\ang{59.21}\sin\ang{20}/\cos\ang{71.34} = 0.9184$, $\beta = \ang{66.66}$।
चूँकि $\cos p = 0.5119 < \sin\ang{40}\times0.9474 = 0.6090$, दूसरी स्थिति
लागू होती है: $\varphi_p = 10 + 180 - 66.66 = \ang{123.34}$\,E। ध्रुव
\ang{71.3}\,N, \ang{123.3}\,E पर है।
\end{solution}

\begin{solution}{exo:b3:palaeomagnetism:7}
(क) $(7.69 - 8.32)/8.32 = \qty{-7.6}{\percent}$ 125 वर्षों में, या
\qty{-6.1}{\percent} प्रति शताब्दी। (ख) $29350/12.6 = 2329$, लगभग 2300 वर्ष।
यह कोई पूर्वानुमान नहीं है: दर स्वयं एक पाँच-वर्षीय मॉडल से
अगले मॉडल तक बदलती है, आघूर्ण अतीत में उलटे बिना ऊँचा और नीचा दोनों रहा है, और
क्षेत्र अनियमित अंतरालों पर उलटता है, जिनका
कोई दर पूर्वानुमान नहीं कर पाती।
\end{solution}

\begin{solution}{exo:b3:palaeomagnetism:8}
20 और \qty{300}{\celsius} के बीच सदिश ने $(6.4, 0.0, 13.6)$ खो दिया: एक
घटक \qty{15.0}{mA/m} का, जिसमें $D = \ang{0.0}$ और
$I = \arctan(13.6/6.4) = \ang{64.8}$। \qty{300}{\celsius} से ऊपर इसने $(25.8, 18.1, 24.6)$ खो दिया:
\qty{40.0}{mA/m}, $D = \arctan(18.1/25.8) = \ang{35.1}$, $I =
\arctan(24.6/31.5) = \ang{38.0}$। दूसरा, जो मूल बिंदु की ओर जाता है, वह
\omterm{def:b3:palaeomagnetism:zijderveld}{अभिलक्षणिक अवशिष्ट चुंबकन} है; पहला, उत्तर की ओर की तीव्र दिशा में, जो
निम्न तापमान पर हटाया गया, वर्तमान क्षेत्र के अनुदिश श्यान अधिलेखन है।
\end{solution}

\begin{solution}{exo:b3:palaeomagnetism:9}
$\lambda = \arctan(\tfrac12\tan(\ang{-35})) = \ang{-19.3}$। $I = \ang{-40}$
और \ang{-30} के लिए: \ang{-22.8} और \ang{-16.1}, अंतराल \ang{6.7}।
$p = \ang{90} - \lambda = \ang{109.3}$, $\cos^2p = 0.1096$:
$\mathrm dp = 2.5\times(1 + 0.329) = \ang{3.3}$ और $\mathrm dm =
5\times\sin\ang{109.3}/\cos\ang{35} = \ang{5.8}$।
\end{solution}

\begin{solution}{exo:b3:palaeomagnetism:10}
(क) $\tan\lambda = \tfrac12$, $\lambda = \ang{26.6}$। $F$ भूमध्य रेखा पर $B_0$ है और
ध्रुवों पर $2B_0$: अनुपात 2। (ख) $\tan I = 2\tan\lambda$ से
$(1 + \tan^2I)\,\mathrm dI = 2(1 + \tan^2\lambda)\,\mathrm d\lambda$ प्राप्त होता है, अतः
\[
  \frac{\mathrm dI}{\mathrm d\lambda} = \frac{2/\cos^2\lambda}{1 + 4\tan^2\lambda}
  = \frac{2}{\cos^2\lambda + 4\sin^2\lambda} = \frac{2}{1 + 3\sin^2\lambda}.
\]
एक त्रुटि $\delta I$ से $\delta\lambda = \tfrac12(1 + 3\sin^2\lambda)\,\delta I$ प्राप्त होता है,
भूमध्य रेखा पर $\tfrac12\delta I$ से ध्रुवों पर $2\,\delta I$ तक: सबसे बड़ी
त्रुटि ध्रुवों के निकट है। (ग) भूमध्य रेखा के निकट नति अक्षांश की तुलना में
दोगुनी तेज़ी से बदलती है, इसलिए नति की दी गई त्रुटि अक्षांश में
केवल आधी त्रुटि उत्पन्न करती है; ध्रुवों के निकट नति अक्षांश के साथ लगभग नहीं बदलती।
\end{solution}

\begin{solution}{exo:b3:palaeomagnetism:11}
(क) $\cos p = \sin(\ang{-35})\sin\ang{70} + \cos\ang{35}\cos\ang{70}\cos\ang{150}
= -0.5390 - 0.2426 = -0.7816$, $p = \ang{141.41}$। फिर $\tan I_x = 2\cot p =
-2.506$, $I_x = \ang{-68.2}$, और $\lambda_x = \ang{90} - p = \ang{-51.4}$।
$\cos D_x = (\sin\ang{70} - \sin(\ang{-35})\cos p)/(\cos\ang{35}\sin p) =
0.9617$, और चूँकि ध्रुव स्थल के याम्योत्तर के पूर्व में है ($\sin\ang{150} >
0$), $D_x = \ang{15.9}$। (ख) $\lambda_o = \arctan(\tfrac12\tan(\ang{-60})) =
\ang{-40.9}$। $R = 30.0 - 15.9 = \ang{14.1}$: खंड का दक्षिणावर्त घूर्णन
ऊर्ध्वाधर अक्ष के परितः। $P = \lambda_x - \lambda_o = -51.4 + 40.9 =
\ang{-10.5}$: संरचना स्थिर प्लेट के अनुसार जितनी होनी चाहिए थी, उससे
\ang{10.5} और उत्तर में स्थित थी, इसलिए खंड तब से लगभग \ang{10.5},
\qty{1170}{km} के बराबर, प्लेट के सापेक्ष दक्षिण की ओर खिसक चुका है।
\end{solution}

\begin{solution}{exo:b3:palaeomagnetism:12}
(क) $\mathrm d\Omega = 2\pi\sin\theta'\,\mathrm d\theta'$ लेने पर,
\[
  \int_0^\theta\frac{\kappa e^{\kappa\cos\theta'}}{4\pi\sinh\kappa}\,2\pi\sin\theta'\,
  \mathrm d\theta' = \frac{e^\kappa - e^{\kappa\cos\theta}}{2\sinh\kappa}
  = \frac{e^\kappa - e^{\kappa\cos\theta}}{e^\kappa - e^{-\kappa}} .
\]
(ख) बड़े $\kappa$ के लिए $e^{-\kappa}$ नगण्य है; $e^\kappa$ से भाग देने पर
शेष रहता है $1 - e^{-\kappa(1 - \cos\theta)} \approx 1 - e^{-\kappa\theta^2/2}$। यह
$1 - e^{-1} = \qty{63}{\percent}$ के बराबर होता है, जब $\kappa\theta^2/2 = 1$, $\theta =
\sqrt{2/\kappa}$ रेडियन, या $\ang{81}/\sqrt\kappa$। (ग) $81/\sqrt{50} =
\ang{11.5}$; सटीक सूत्र \qty{63.1}{\percent} देता है।
\end{solution}

\begin{solution}{pb:b3:palaeomagnetism:1}
\textbf{1.} लाल संस्तर में हेमेटाइट, लौह ऑक्साइड और हाइड्रॉक्साइड से,
अवसाद जमने के काफ़ी बाद बन सकता है, और तब वह उस बाद के काल का \omterm{def:b3:palaeomagnetism:nrm}{रासायनिक अवशिष्ट चुंबकन} देता है।
यदि संस्तर जमने के तुरंत बाद मोड़े गए थे, तो धनात्मक वलन-परीक्षण,
या पास-पास स्थित दोनों ध्रुवताओं के संस्तरों पर धनात्मक उत्क्रमण-परीक्षण,
यह दिखा देगा कि अवशिष्ट चुंबकन प्रारंभिक है।

\textbf{2.} $(\cos\ang{37}\cos\ang{50}, \cos\ang{37}\sin\ang{50}, \sin\ang{37})
= (0.5134, 0.6118, 0.6018)$।

\textbf{3.} $R = \sqrt{3.8943^2 + 2.6500^2 + 3.5717^2} = 5.9114$। $D =
\arctan(2.6500/3.8943) = \ang{34.2}$, $I = \arcsin(3.5717/5.9114) = \ang{37.2}$।

\textbf{4.} $k = 5/(6 - 5.9114) = 56.5$। $20^{1/5} - 1 = 0.8206$ लेने पर,
$\cos\alpha_{95} = 1 - (0.0886/5.9114)\times0.8206 = 0.98771$, $\alpha_{95} =
\ang{9.0}$।

\textbf{5.} B: $k = 6/(7 - 6.8979) = 58.8$; $20^{1/6} - 1 = 0.6475$,
$\cos\alpha_{95} = 1 - (0.1021/6.8979)\times0.6475 = 0.99042$, $\alpha_{95} =
\ang{7.9}$। C: $k = 7/(8 - 7.8707) = 54.1$; $20^{1/7} - 1 = 0.5340$,
$\cos\alpha_{95} = 1 - (0.1293/7.8707)\times0.5340 = 0.99123$, $\alpha_{95} =
\ang{7.6}$।

\textbf{6.} B का $k$ सबसे बड़ा है, अर्थात् सबसे कसा समूह; C का
$\alpha_{95}$ सबसे छोटा है, इसलिए उसका माध्य सबसे सुनिश्चित है, क्योंकि उसके सबसे अधिक स्थल हैं।
तीनों एक-दूसरे के निकट हैं: तीनों संरचनाएँ भली-भाँति समूहित हैं।

\textbf{7.} $\lambda_A = \arctan(\tfrac12\tan\ang{37.17}) = \ang{20.76}$\,N,
$\lambda_B = \ang{32.15}$\,N, $\lambda_C = \ang{45.33}$\,N।

\textbf{8.} $p = 90 - 20.76 = \ang{69.24}$। $\sin\lambda_p =
\sin\ang{47}\cos\ang{69.24} + \cos\ang{47}\sin\ang{69.24}\cos\ang{34.23} =
0.2593 + 0.5272 = 0.7865$: $\lambda_p = \ang{51.86}$\,N।

\textbf{9.} $\sin\beta = \sin\ang{69.24}\sin\ang{34.23}/\cos\ang{51.86} =
0.8516$, $\beta = \ang{58.40}$। चूँकि $\cos p = 0.3545 < \sin\ang{47}\times
0.7865 = 0.5752$, इसलिए दूसरी स्थिति लागू होती है, जहाँ $\cos\beta < 0$:
$\varphi_p = 15.00 + 180 - 58.40 = \ang{136.60}$\,E।

\textbf{10.} $\cos^2p = 0.1257$: $\mathrm dp = 4.50\times1.377 = \ang{6.2}$;
$\mathrm dm = 8.99\times\sin\ang{69.24}/\cos\ang{37.17} = \ang{10.5}$।

\textbf{11.} \ang{51.9}\,N (\qty{60}{Ma}) से \ang{66.5}\,N (\qty{30}{Ma}) होते हुए
\ang{87.1}\,N (\qty{4}{Ma}) तक, लगभग एक ही याम्योत्तर पर
(\ang{136}--\ang{139}\,E): भौगोलिक ध्रुव की ओर एक पथ, जिसे
सबसे युवा संरचना अपने शंकु के भीतर पहुँचती है।

\textbf{12.} \qty{22}{m} पर स्थित राख संस्तर दूसरे सामान्य क्षेत्र में पड़ता है,
19.7--\qty{24.6}{m}; \qty{3.2}{Ma} के आसपास की आयु पर वह क्षेत्र C2An.2n
(3.207--3.116)। फिर 4.2--\qty{14.8}{m} का क्षेत्र C2An.3n है, और
29.2--\qty{54.0}{m} का क्षेत्र C2An.1n है, उत्क्रमित आधार C2An के नीचे वाले ध्रुवता क्रॉन का है,
और उत्क्रमित शीर्ष उसके ऊपर वाले ध्रुवता क्रॉन का।

\textbf{13.} $(19.7 - 4.2)/(3.596 - 3.207) = 15.5/0.389 = \qty{39.8}{m/Myr}$;
$(54.0 - 24.6)/(3.116 - 2.581) = 29.4/0.535 = \qty{55.0}{m/Myr}$।

\textbf{14.} आधार: $3.596 + 4.2/39.8 = \qty{3.70}{Ma}$। शीर्ष: $2.581 -
10.0/55.0 = \qty{2.40}{Ma}$।

\textbf{15.} केवल उत्क्रमण सीमाओं का काल-निर्धारण हुआ है, जो कुछ सौ हज़ार वर्ष के अंतर पर हैं;
उनके बीच दर स्थिर मानी गई है। अवधसान, अवसाद आपूर्ति, या जलवायु में
परिवर्तन, या निक्षेपण में अंतराल, उस दर को बदल देगा, और
एक विराम छोटे क्षेत्र को भी हटा सकता है और मिलान को भ्रमित कर सकता है।

\textbf{16.} $\cos p = \sin\ang{47}\sin\ang{78} + \cos\ang{47}\cos\ang{78}
\cos\ang{175} = 0.7154 - 0.1413 = 0.5741$, $p = \ang{54.96}$। $\tan I_x =
2\cot p = 1.4023$, $I_x = \ang{54.5}$; $\lambda_x = \ang{35.04}$\,N।

\textbf{17.} $\cos D_x = (\sin\ang{78} - \sin\ang{47}\times0.5741)/
(\cos\ang{47}\sin\ang{54.96}) = 0.99976$; ध्रुव स्थल के
याम्योत्तर के पूर्व में है ($\sin\ang{175} > 0$), इसलिए $D_x = \ang{1.3}$।

\textbf{18.} $I_x$ = \ang{54.5}, \ang{60.0} और
\ang{63.6} से अपेक्षित \omterm{def:b3:palaeomagnetism:vgp}{पुराअक्षांश}: \ang{35.03}, \ang{40.89} और \ang{45.21}। $P_A = 35.03 - 20.76 =
\ang{14.3}$, $P_B = 40.89 - 32.15 = \ang{8.7}$, $P_C = 45.21 - 45.33 =
\ang{-0.1}$।

\textbf{19.} $R_A = 34.2 - 1.3 = \ang{32.9}$, $R_B = 23.9 - 2.9 = \ang{21.0}$,
$R_C = 3.4 - 2.5 = \ang{0.9}$, सभी दक्षिणावर्त।

\textbf{20.} $\Delta D_o = \arcsin(\sin\ang{8.99}/\cos\ang{37.17}) = \ang{11.3}$;
$\Delta D_x = \arcsin(\sin\ang{3.0}/\cos\ang{54.5}) = \ang{5.2}$; $\Delta R_A =
\sqrt{11.3^2 + 5.2^2} = \ang{12.4}$।

\textbf{21.} $R_A = \ang{32.9} \pm \ang{12.4}$ और $R_B = \ang{21.0} \pm
\ang{13.8}$ सार्थक हैं; $R_C = \ang{0.9} \pm \ang{17.9}$ नहीं है।

\textbf{22.} साठ मिलियन वर्ष पहले केस्टा लगभग \ang{14} अक्षांश
दक्षिण में था, जहाँ स्थिर महाद्वीप उसे रखता है, और अपने वर्तमान अभिविन्यास से लगभग \ang{33}
वामावर्त घूमा हुआ था। \qty{30}{Ma} तक वह उन \ang{14} में से लगभग
\ang{6} उत्तर की ओर आ गया था, और उन \ang{33} में से \ang{12} दक्षिणावर्त घूमा
था; \qty{4}{Ma} तक वह महाद्वीप पर अपने वर्तमान स्थान पर पहुँच चुका था,
घूर्णन शेष नहीं रहा। वह महाद्वीप की सीमा के साथ उत्तर की ओर खिसका,
जाते हुए दक्षिणावर्त घूमता रहा, और \qty{4}{Ma} तक अपनी जगह पर था।

\textbf{23.} $47.00 - 20.76 = \ang{26.24}$, या $26.24\times111.19 =
\qty{2918}{km}$। इसमें स्थिर महाद्वीप की अपनी गति भी शामिल है;
उसके सापेक्ष केस्टा $P_A = \ang{14.3}$ खिसका।

\textbf{24.} $I = \ang{28.18}$ और \ang{46.16} के लिए $\lambda = \ang{15.00}$ और
\ang{27.50} प्राप्त होते हैं: अर्ध-चौड़ाई \ang{6.25} है, और एक मानक विचलन
$6.25/1.96 = \ang{3.19}$ है, या \qty{355}{km}। संरचना A के काल से केस्टा
उत्तर की ओर
\textbf{\qty[text-series-to-math=true]{2918 \pm 355}{km}} खिसका है।
\end{solution}
